% SPDX-License-Identifier: CC-BY-4.0
%
% Preprint source (LaTeX). Primary venue: IACR Cryptology ePrint Archive
% (eprint.iacr.org; category: Cryptographic protocols / Applications; no
% endorsement required). Fallback: Zenodo (DOI). arXiv cs.CR is locked for this
% topic for us; see papers/preprint/STAGING.md for the submission checklist.
%
% Content is a faithful carriage of papers/authorization-receipts-preprint.md,
% with evidence-backed claims and reference revisions synchronized to repo state.
% Canonical Markdown SHA-256: 90de77af844d5c1bfbd15645d095bc01a8b78d3a8d4c5bb15ba09553ff99c151
% Every quantitative claim is traceable to a repository artifact:
%   - conformance 21 suites / 331 vectors -> conformance/conformance-manifest.json
%     (totals.suites=21, totals.vectors=331)
%     manifest_sha256=7614f20f281c87f688f689a3246ca4ea5d2e76dc1b36e525e783b4eb0e60bc87
%   - all Tamarin lemma output blocks -> formal/PROOF_STATUS.md (verbatim)
%   - TLA+ 413,137 states / 45,342 distinct / 26 invariants -> formal/PROOF_STATUS.md
%   - Alloy 15 + 7 + 6 + 4 = 32 assertions -> formal/PROOF_STATUS.md
%   - Rust external verifier commit 7faba360..., 164 vectors, 353+6 hostility
%     -> conformance/external/rust-cleanroom-jdieselny.v1.json
%
\documentclass[11pt]{article}

\usepackage[T1]{fontenc}
\usepackage[utf8]{inputenc}
\usepackage{lmodern}
\usepackage[margin=1in]{geometry}
\usepackage{amsmath,amssymb}
\usepackage{textcomp}
\usepackage{microtype}
\usepackage{enumitem}
\usepackage{fancyvrb}
\usepackage{parskip}
\usepackage[hidelinks]{hyperref}
\usepackage{xurl}
\urlstyle{same}

\hypersetup{
  pdftitle={Action-Bound Injective Authorization under Signer Harvesting: A Symbolic Security Analysis},
  pdfauthor={Iman Schrock (EMILIA Protocol, Inc.)},
  pdfsubject={cs.CR},
  pdfkeywords={action-bound authorization, injective authorization, signer harvesting, symbolic analysis, Tamarin, Dolev-Yao, human approval evidence}
}

% Inline monospaced code with safe underscores/hyphens.
\newcommand{\code}[1]{\texttt{#1}}
% Long, unbreakable identifiers (hex digests, long draft names): allow wrapping.
\usepackage{seqsplit}
\newcommand{\hexid}[1]{\texttt{\seqsplit{#1}}}

% Verbatim style for machine-checked tool output.
\DefineVerbatimEnvironment{tool}{Verbatim}{fontsize=\footnotesize,frame=single,framesep=4pt}

\title{Action-Bound Injective Authorization under Signer Harvesting: A Symbolic
Security Analysis}

\author{Iman Schrock\\
EMILIA Protocol, Inc.\\
\texttt{team@emiliaprotocol.ai}\\
ORCID: \href{https://orcid.org/0009-0004-0290-5433}{0009-0004-0290-5433}}

\date{Preprint. August 2026.}

\begin{document}
\maketitle

\begin{center}
\emph{The construction analyzed here is also specified, in engineering form, as a
set of individual Internet-Drafts (\code{draft-schrock-ep-authorization-receipts-10}
and companions). Those are individual submissions with no working-group status;
this paper is self-contained and does not depend on them.}
\end{center}

\begin{abstract}
Digital signatures authenticate messages, but signature authenticity alone does
not imply that an authorization is non-transplantable to another action or
injectively consumable at an enforcement point. We study this
protocol-composition problem when the service that renders requests, routes
signoffs, and retains records is untrusted. We define \emph{Action-Bound
Injective Authorization (ABIA)}: an approver-held key signs an injectively
encoded authorization context naming one exact action, and a relying party
accepts only after verifying that context and atomically consuming its
authorization within a stated domain. The adversary controls the network and
may obtain valid signatures on contexts of its choice, which models signer
harvesting by a malicious operator, but does not compromise the approver keys.
In the symbolic trace model, the receipt core establishes action
non-transplantability and user-verification-gated authenticity; the consumption
rule establishes injective terminal acceptance within one shared domain; and the
2-of-2 quorum model establishes signer distinctness and initiator exclusion when
those values are signed. Removing consumption, initiator binding, or exact
registry-view binding yields concrete traces while leaving the signature
primitive ideal. Four Tamarin models verify 17 all-traces properties and three
executable witnesses; eight deliberately weakened comparisons are falsified.
The result is a protocol-security theorem under Dolev--Yao assumptions, not a
new signature primitive or a computational reduction. TLA+, Alloy,
implementation, and conformance artifacts validate the surrounding construction
but do not replace the symbolic result. The analysis does not prove
natural-person identity, presentation fidelity, policy correctness,
log completeness, or downstream effects.
\end{abstract}

\section{Introduction}

\subsection{The testimony-versus-evidence problem}

Consider a dispute that is now routine in outline and will shortly be routine in
fact. An AI agent operated by a financial-services firm releases a wire transfer.
The transfer is later challenged, and the firm produces its approval records: a
workflow-tool entry showing that a user clicked ``approve'' at a certain time.

Every element of that record is an assertion by the firm. The database holding
the approval row is writable by the firm's administrators; the mapping between
the row and the executed action is maintained by the firm's software; the
timestamp is the firm's clock. An auditor, counterparty, insurer, or court that
wants to rely on the record must trust the operator of the approval system, and
the operator is frequently the party whose conduct is in question. The record is
testimony.

Existing authorization infrastructure does not change this, because it answers a
different question. Identity and access management establishes that an actor is
authenticated and authorized \emph{in general}: it grants sessions and scopes,
not decisions about individual actions. Fraud that occurs inside a valid session
through approved channels, such as a business-email-compromise-driven beneficiary
change, is invisible to session-level controls. Where human approval does exist,
it is a click recorded in a mutable store. Three gaps follow, and they are
structural rather than implementation accidents:

\begin{enumerate}[leftmargin=*]
\item \textbf{The action gap.} Authorization is granted to sessions and scopes,
not to individual actions with concrete parameters.
\item \textbf{The accountability gap.} Typical approval workflows provide no
independent cryptographic binding from an enrolled approver key to the canonical
bytes of a specific action; the approval record is producible and alterable by the
operator.
\item \textbf{The verification gap.} Those workflow records require third parties
to trust the operator's logs rather than carrying a portable artifact that can be
checked under independently selected trust roots.
\end{enumerate}

The authorization receipt is a narrow construction aimed at exactly these gaps:
before an irreversible action executes, an enrolled approver key signs a canonical
context binding the exact action; the authorization reaches a terminal state at
most once within the executor's shared atomic consumption domain; and the
resulting receipt is verifiable offline for as long as its algorithms and
independently selected trust material remain acceptable or are renewed.

\subsection{The security question and boundary}

The claim is deliberately limited. A receipt is evidence that a key enrolled
under an approver identifier signed a canonical context binding one action, a
committed policy version, and a validity window, and that the authorization
reached a terminal state once within a shared atomic consumption domain. It is not
evidence that the enrolled key holder is a particular natural person, that the
action was a good idea, that the approver understood the business context, that an
independent executor did not accept a replay, or that the surrounding deployment
cannot be bypassed. Verification proves signature, binding, and log-inclusion
integrity; it never proves business correctness or a downstream effect. The
organization publishing this work is not an auditor, regulator, or insurer; the
artifacts described here support the judgments of such parties and conclude
nothing on their own.

\subsection{What standard signature security does not provide}

The contribution is not the claim that a signature can authenticate a byte
string. That is the primitive assumption. The cryptographic question is what
must be added around that primitive before a relying party may treat a signed
authorization as evidence for one exact action and one terminal admission.

Let $\mathsf{Sign}_H(x)$ be an honest approver's signing operation on
authorization context $x$, and let $\mathsf{Accept}_D(x)$ be terminal acceptance
in atomic domain $D$. The signer-harvesting adversary may choose $x$ and obtain
$\mathsf{Sign}_H(x)$ for arbitrary actions, then replay, splice, or reorder the
resulting artifacts. It wins one of three protocol games if:

\begin{enumerate}[leftmargin=*]
\item \textbf{Transplantation:} $\mathsf{Accept}_D(x)$ occurs without an
honest signature on the same canonical context $x$ (an approval for one action
is used for another).
\item \textbf{Non-injective acceptance:} two terminal acceptances in the same
$D$ consume the same authorization context (the same valid receipt is replayed).
\item \textbf{Quorum substitution:} a quorum commits an action without distinct
enrolled approvers signing the same action and initiator, or allows the initiator
to occupy an approver slot.
\end{enumerate}

For the 2-of-2 profile analyzed here, we call the conjunction of these
properties \emph{ABIA security}. This is a
protocol-level security notion, not a replacement for standard signature
security: a scheme may remain perfectly unforgeable while a verifier that omits
consumption or leaves the initiator outside the signed context loses ABIA. The
negative traces in Section 7 are therefore separation results. They show that
the missing checks are not editorial detail and that a valid signature, a log
inclusion proof, or a quorum count does not supply the omitted protocol
property by itself.

\subsection{Contributions}

\begin{itemize}[leftmargin=*]
\item \textbf{A protocol abstraction and ABIA security definition.} We isolate
action-bound authorization from its agent and workflow integrations, define the
signer-harvesting games, terminal acceptance relation, shared consumption domain,
quorum condition, and relying-party key pinning, and state the exact properties
a relying party is entitled to infer.
\item \textbf{A symbolic security theorem and separations.} Four Tamarin models
derive the receipt and quorum properties rather than assuming that an accepted
signature is enough. The strict models verify 17 all-traces properties and three
executable witnesses; weakened models produce concrete replay, initiator-binding,
and stale-registry-view attacks while the signature algebra remains ideal.
\item \textbf{An operator-untrusted adversary model.} The party that renders the
action, routes approval, and keeps the record may drive an honest approver to
sign any action chosen by the attacker. This removes the honest-renderer
assumption and makes the result about protocol evidence rather than ordinary
session authorization.
\item \textbf{A precise construction boundary.} Canonical action bytes,
approver-held keys, atomic terminal consumption, and offline verification are the
analyzed core. A host-record binding profile and evidence-sufficiency graph are
secondary compositions and are not claimed as new cryptographic primitives.
\item \textbf{Bounded validation.} The reference implementation,
21-suite/331-vector conformance battery, and separately authored Rust verifier
validate encoding and interoperability. They are not presented as independent
security proofs.
\end{itemize}

The paper does not claim a new signature, hash, or transparency-log primitive.
Its cryptology contribution is the ABIA protocol result: under a
signer-harvesting adversary, exact action binding and terminal injectivity are
separate obligations, the state and context bindings needed for them can be
machine-checked, and removing those bindings produces attacks without forging
the underlying signatures.

\subsection{The result in one statement}

\textbf{Theorem 1 (symbolic ABIA result).} Let $x=(a,p,i,n,w,v)$ be the
injectively encoded authorization context, where $a$ is the canonical action,
$p$ the policy commitment, $i$ the initiator, $n$ the one-time nonce, $w$ the
validity window, and $v$ the approver slot. Let $D$ be the executor's shared
atomic consumption domain. In the trace semantics of the models in Section 7,
with pinned key material, uncompromised signing keys, and the stated injective
encoding assumption:

\begin{itemize}[leftmargin=*]
\item \textbf{Receipt core:} an accepted receipt implies a prior
user-verification-gated signature over the exact $(a,n)$ context; an honest
signature over another action cannot be transplanted.
\item \textbf{Terminal injectivity:} within $D$, two terminal acceptances cannot
consume the same $(H(x),n)$ authorization. This guarantee is not portable across
independent executors that do not share $D$.
\item \textbf{Quorum integrity:} the 2-of-2 model requires two distinct enrolled
approver identities, each signing the same action and initiator; the initiator
cannot fill an approver slot.
\end{itemize}

The counterexamples are part of the theorem's boundary. Removing the consumption
check admits a replay trace; removing the initiator from the signed context admits
a quorum/separation-of-duties trace; and omitting the exact pinned registry view
admits a stale-view trace. These are protocol failures, not failures of the
underlying signature primitive. The composed models extend the same reasoning to
challenge registration, authority scope, reliance configuration, and action-keyed
execution; they do not silently promote those additional fields into the focused
receipt theorem.

\textbf{Proof sketch.} In \code{ep\_receipt\_core.spthy}, the only rule that
emits an honest signature consumes a linear user-verification fact for the same
action and nonce. The acceptance rules can fire only when Tamarin's signature
equation holds under the pinned key, so the equality restriction yields the
exact signed context. The \code{ConsumeOnce} restriction makes two checked
acceptances of the same authorization impossible. In
\code{ep\_quorum\_core.spthy}, the commit rule requires two verification
equations under distinct pinned keys and explicitly excludes the initiator;
removing the initiator field from the signed term makes the corresponding
separation-of-duties query falsifiable. The composed models apply the same
derivation to the pinned registry view and challenge/action key. Thus the
positive result and the negative traces use the same symbolic signature theory;
only the protocol rules and restrictions differ.

\section{Threat Model}

Three adversarial settings drive the design. They are treated together because a
mechanism that addresses one while silently assuming away another produces
evidence that fails exactly when it matters.

The adversary here is stronger than the one standard authorization models carry,
and the difference is the point of the work. Session and delegation systems, from
OAuth scopes and Rich Authorization Requests through transaction tokens and
delegation receipts, reason about which principal may act, and they assume that an
honest principal signs only what a correct protocol run presents to it. This
setting removes that assumption. The operator that renders the action to the
approver, routes the approval, and later holds the record is untrusted for
evidence purposes, and it can present an honest human with any action it chooses
to have signed. The question is therefore not whether the actor was authorized in
general, which the delegation layer answers, but whether an accepted record can
still prove that an uncompromised key enrolled to the cited approver identifier
signed a context binding exactly the accepted action when the party producing the
record is the one whose conduct is later in dispute. The symbolic analysis in
Section 7 is run against exactly this adversary, which is what makes its
authenticity and no-replay results say something the delegation layer cannot.

\subsection{Compromised or malicious operator}

The party running the orchestration service (policy registry, signoff routing,
log) is not trusted for evidence purposes. Under the recommended key custody
classes and the model's uncompromised-key assumption (Section 4.4), a compromised
operator can deny service and can fail to route signoff requests; it cannot forge
a signature under an enrolled approver key it does not hold. Within a shared
atomic consumption domain, a second presentation of a genuine authorization is
rejected. Independent executors require a shared consumption service or equivalent
coordination to obtain that replay guarantee.

Two operator-compromise paths remain open and are part of the model rather than
claimed away. First, an operator that controls the signing client's rendering can
harvest a \emph{genuine} signature over an action the approver misunderstood; this
is the presentation-attack family of Section 10.2, and it is why independently
authored rendering surfaces are required for high-value policies. Second, an
operator that unilaterally controls the directory of enrolled approver keys can
enroll a key it controls under a legitimate approver's name, relocating the
forgery rather than preventing it; Section 4.4 describes the directory-authority
controls. The accurate statement of the modeled property is therefore: under the
uncompromised-key assumption, the operator cannot forge a signature under an
enrolled approver key. The stronger statement, that the operator cannot obtain an
unauthorized approval at all, additionally requires the directory-authority and
independent-rendering controls to be deployed.

\subsection{Compromised agent}

The initiating agent is identified but never trusted. A prompt-injected or
otherwise compromised agent can propose any action, state any escalation reason,
and attribute itself to any external identity. The design consequence is uniform:
everything the initiator supplies is treated as a claim. Injection can change what
the initiator \emph{proposes}; it cannot change what a human \emph{approves} on
the human's own hardware, because the device-bound signature is produced outside
the model's context. The protocol additionally enforces separation of duties: an
initiator can never occupy an approver slot for its own action. This separation is
one of the properties the symbolic quorum model verifies (Section 7).

\subsection{Post-hoc dispute}

After execution, the parties evaluating the evidence (auditors, counterparties,
insurers, courts, regulators) may distrust every online service involved, may be
examining events years old, and may have no network path to the original
operator, which may no longer exist. The design consequence is offline
verifiability: a receipt must be checkable using only its own bytes, pinned or
directory-proven approver key material, and a log checkpoint carried inside the
receipt. What offline verification can and cannot establish in this setting is
bounded precisely in Sections 4.3 and 10.3.

\subsection{Out of scope}

The protocol binds approvals to \emph{approver identifiers} whose keys are
enrolled in a directory; proving that the holder of an identifier is a particular
natural person is the job of the identity-proofing layer that populates the
directory, not of the receipt format. Collusion among distinct enrolled humans,
one human controlling multiple enrolled identities, and coercion of an approver
are likewise not defeated by the mechanism; receipts make such events
attributable, which raises their cost, and no more than that. The symbolic
analysis is explicit about this boundary: it proves the distinct-identity and
no-self-approval property, and nothing about whether two distinct identities are
two independent wills (Section 7.4).

\section{The Receipt Construction}

\subsection{Actions as canonical bytes}

An action is a single proposed operation with concrete parameters: one wire
transfer to one beneficiary for one amount, one credential rotation of one key. It
is represented as a JSON Action Object containing at minimum the action type,
target system and resource, parameters, initiator identity, policy reference, and
request time. The object is serialized under the JSON Canonicalization Scheme
(JCS, RFC 8785), and the \emph{action hash} is the SHA-256 digest of the canonical
serialization. Implementations must reject an approval request whose action hash
does not match a locally recomputed hash of the presented object. Sensitive
parameter values may be carried as salted hashes, provided the executing system
can recompute them; the binding property is preserved because the hash commits to
the committed values.

Canonical bytes are the foundation of every later property: an approval bound to a
description or a workflow-ticket identifier binds to something the operator can
re-narrate; an approval bound to the SHA-256 of canonical bytes binds to exactly
one action. The symbolic models of Section 7 make this abstract but load-bearing:
they treat message encoding as injective (the standard symbolic assumption) and
derive no-replay-across-actions from the fact that each signature is over the
action's hash, so the attacker cannot repurpose a signature over one action term
as a signature over another.

\subsection{The authorization context and the signoff}

For each required approver, the orchestrator constructs an Authorization Context: a
canonical JSON structure containing the action hash, the policy identifier
\emph{and} a policy hash committing to the exact policy version evaluated, the
initiator identity, the approver identifier and slot index, the required approval
count, a one-time nonce of at least 128 bits of CSPRNG output, issuance and expiry
times, and a hash chaining the attempt to the issuing log's most recent receipt.
The context is JCS-canonicalized; the \emph{context hash} is its SHA-256 digest;
the approver signs the context hash.

Three bindings deserve emphasis, and each is exactly the kind of load-bearing
detail a symbolic prover exposes. First, the policy commitment is exact: a
signature over a context with policy hash X must not satisfy a requirement
evaluated under policy hash Y, even for the same policy identifier. Approvals
cannot migrate across policy versions. Second, the initiator identity is bound
\emph{inside} the signed context, not merely tracked as an executor-local label;
Section 7.3 records that an earlier quorum model that omitted this binding was
falsified, because separation of duties then held only at the executor and not in
the signatures themselves. Third, the approver must be shown, at signing time, a
faithful human-readable rendering of the Action Object, not only its hash;
interfaces that display a different action than the one hashed constitute a
presentation attack (Section 10.2).

The context may optionally carry two further members, both claims by parties the
protocol identifies but never trusts. An \emph{initiator attestation} records the
initiator's own stated reason for escalating to a human, as a closed enumeration
(irreversibility, magnitude, uncertainty, novelty, authority gap, policy rule)
plus an optional rule identifier and a length-capped free-text statement. An
\emph{agent binding} attributes the action to an external agent identity and,
optionally, the external delegation under which the agent acted, with an optional
observed-at timestamp against which a relying party may enforce evidence
freshness, fail-closed. Because JCS serializes every member present, both are
covered by the approver's signature: the receipt proves the stated reason and
claimed attribution were part of what the approver signed. Neither is proof of the
initiator's internal state, the agent's identity, or the delegation's validity,
and policy engines must not use attestation content to relax any threshold: the
initiator must gain nothing by saying the right words.

The signoff itself carries the context hash, the signature, the key custody class,
the approver key identifier, and, for device-bound keys, the WebAuthn assertion
material. For device-bound keys the WebAuthn challenge must equal the context hash
and the authenticator's user-verification flag must be set, so the signature
attests both possession of the enrolled authenticator and a local user-verification
event. In the symbolic models this user-verification requirement is modeled
structurally as a linear precondition that only a user-verification step can
satisfy, so no trace contains a signature without a preceding user-verification
event over exactly that action and nonce (Section 7.2); the model assumes the
authenticator enforces user verification and does not itself prove WebAuthn
internals. Denials are signed over the same context hash with a denial envelope,
so refusals are equally attributable to the enrolled key and equally terminal:
``the enrolled approver key signed a refusal of this action'' is a positive,
verifiable fact, not an absence.

\subsection{Core protocol interface and adversarial queries}

The cryptographic core has four roles: an approver key holder, an operator that
may be malicious, a relying party, and an enforcement domain $D$. The operator
constructs a canonical context $x$, obtains an approver signature, and returns
the resulting artifact to the network. The relying party runs
$\mathsf{Verify}(x,sig,K)$ against independently pinned key material, then runs
the atomic operation $\mathsf{Consume}_D(H(x),n)$ before terminally admitting the
action. A receipt may include a transparency-log inclusion proof, but that proof
authenticates publication; it does not replace either the signature check or the
consume operation.

The formal adversary is allowed to call the honest signing rule on arbitrary
actions and initiators. This models an operator that can make a genuine signing
request look legitimate to the signing client. The models do not grant the
adversary the signing key: $\mathsf{RevealLtk}$ is an explicit compromise event,
and each authenticity lemma is conditioned on its absence before acceptance.
This separation is important. The theorem addresses misuse of valid signatures;
it does not claim resilience after key extraction or against a compromised
trusted display path.

\subsection{One-time consumption as enforcement-point state}

An authorization attempt moves through a small state machine:

\begin{tool}
REQUESTED -> {PARTIALLY_APPROVED}* -> APPROVED -> COMMITTED
          \-> DENIED                          \-> EXPIRED
          \-> EXPIRED
\end{tool}

COMMITTED, DENIED, and EXPIRED are terminal. The construction's invariants,
maintained as machine-checked models (Section 7) and required of conforming
implementations, are:

\begin{itemize}[leftmargin=*]
\item \textbf{ConsumeOnce.} Within one shared atomic consumption domain, a nonce
transitions to a terminal state at most once; any second presentation to that
domain is rejected as a replay.
\item \textbf{BindingMatch.} A signoff satisfies only the context, and therefore
only the action hash, that it signs.
\item \textbf{TerminalIrreversibility.} No transition exits a terminal state.
\item \textbf{SelfApprovalImpossible.} For every signoff, the approver differs
from the initiator; under multi-approver policies, approvers are pairwise distinct
and each distinct from the initiator.
\item \textbf{NoBypassWrite.} A COMMITTED state is reachable only through the full
sequence, and a conforming verifying executor does not execute without verifying
it.
\end{itemize}

The consumption record is enforcement-point state, not log decoration, and its
necessity is not merely asserted: the symbolic core model (Section 7.2) removes the
consumption check and the prover immediately falsifies injective acceptance,
producing a same-receipt replay trace with no forgery and no key compromise.
Adding the check back restores injectivity. Within the shared atomic consumption
domain, an authorization, once consumed or refused, is terminally unusable; a
later presentation to that domain is detectable and must be rejected. Independent
executors obtain the same guarantee only if they coordinate through that domain or
an equivalent action-keyed atomic admission service. This separates an admitted
authorization from a reusable bearer credential, but it does not by itself prove
that a downstream external system produced exactly one effect.

Deployment topology is graded honestly rather than assumed. In the strongest
conformance class, the system of record itself (payment switch, registry,
deployment controller) verifies the authorization bundle before executing and
refuses otherwise. A middleware class intercepts between the agent and the
executing credential, which defends against agent error and prompt injection but
can be bypassed by an operator with code control. An evidence-only class makes no
enforcement claim at all. Implementations must declare their class in the receipts
they produce and must not state a stronger class than deployed; the distance
between ``the construction has these properties'' and ``your deployment is
unbypassable'' is the most common overclaim in this category.

\subsection{The trust receipt and offline verification}

Upon commitment, the orchestrator assembles the Trust Receipt: the full Action
Object and its hash, every authorization context, every signoff, the consumption
record, a Merkle inclusion proof for the receipt leaf against a signed log
checkpoint (tree size, root hash, log signature, log key identifier), and
inclusion proofs for the approver key entries in the approver directory.
Operator-produced commitments and receipt-log material are signed with Ed25519;
approver signoffs under the device-bound class are ES256 (P-256) or Ed25519 where
the authenticator supports it.

A verifier holding the receipt, a trusted log public key, and a trusted directory
root or pinned approver keys, with \emph{no network access}, establishes six
things: (1) the action hash recomputes from the canonical Action Object; (2) each
context hash recomputes and commits to the action hash, the policy hash, and a
distinct approver; (3) each signature, including the WebAuthn assertion where
present, verifies over the context hash against a key entry whose validity window
contains the issuance time; (4) separation of duties holds and the approval count
satisfies the policy; (5) the Merkle inclusion proof verifies against the
checkpoint root and the checkpoint signature verifies against the log key; (6)
signing and commitment times fall within the validity window.

Step (5) is what distinguishes this design from log-access designs: the checkpoint
travels inside the receipt, so verification requires no query to the log.

One further property is deliberate and load-bearing: no step of offline
verification relies on a message authentication code, a shared secret, or any
symmetric primitive. A symmetric construction, such as an HMAC-chained audit log,
is verifiable only by a party holding the same secret as the producer, precisely
the party the evidence is meant to constrain; its keeper can rewrite its own
history undetectably. Every artifact on the verification path is bound by an
asymmetric signature whose verifying key the relying party holds independently of
the issuer.

Equally important is what offline verification does \emph{not} establish, and the
specification requires implementations to say so. It establishes authenticity as
of commit time, not currency: a receipt whose approver key was revoked an hour
after commitment still verifies, because the artifact is evidence of validity at
commit time. And it establishes inclusion in \emph{a} tree whose head the log
operator signed, not that this tree is the only tree the operator has shown the
world; split-view detection requires online activity such as gossip or witness
cosigning. A relying party with freshness or revocation requirements must
additionally consult a current directory head and checkpoint online.
Implementations must not describe offline verification as establishing that a
receipt is ``currently valid.''

\subsection{Key custody and the approver directory}

Key classes classify custody of the approver's signing key and nothing else. Class
A keys are generated and held in a platform authenticator or security key and
exercised via WebAuthn with user verification required; this is the recommended
class and the one the symbolic models take as their setting. Class B keys are
software keys in the approver's client environment, acceptable where WebAuthn is
impractical, with the reduced assurance noted in receipts. Class C, in which the
operator signs on the approver's behalf after authenticating them, exists only to
describe pre-existing deployments; receipts produced under it must be labeled as
such, and relying parties should treat them as operator assertion, not approver
signature.

Approver public keys are enrolled in a signed, Merkle-tree-structured approver
directory maintained per organization, and a receipt carries an inclusion proof of
the relevant key entry, so verification needs no live directory access. Directory
authority is itself a trust root and must not default to the operator: the
directory tree head must be signed by an organization-controlled key, and where
directory \emph{operation} is delegated to the operator, every enrollment must
carry a second-party attestation by an organization administrator key or a quorum
of already-enrolled approvers. A verifier presented with a directory head signed
only by an operator-held key treats the whole structure as operator assertion,
equivalent in assurance to Class C, regardless of the custody class of individual
signoffs. The rationale is the relocation argument of Section 2.1: an operator
that cannot forge signatures but can decide which keys count has not been removed
from the trust path. The symbolic models abstract this directory as a single
out-of-band key-pinning step; modeling the directory protocol and its trust root
is stated future work (Section 7.4), not a proven property.

The directory is also the explicit slot where identifier-to-person binding lives.
The receipt format proves that a key enrolled under a given identifier signed the
exact action; the strength of the claim that the identifier names a particular
human is a property of the directory authority and whatever identity-proofing
layer feeds it, and it is deliberately not smuggled into the receipt's own claims.

\section{Quorum and Distinctness}

High-consequence actions frequently require more than one approver, and
multi-party approval introduces a failure mode that single-approver analysis
misses: counting signatures instead of humans. This is the section the symbolic
quorum model of Section 7.3 analyzes directly.

Under an m-of-n policy, each required approver receives and signs an individual
authorization context sharing the same action hash and nonce family but carrying a
distinct approver index. Commitment occurs only when the required number of valid,
distinct signoffs exists before expiry. Partial approval confers no authority
whatsoever: a verifying executor presented with fewer than the required signoffs
refuses. The symbolic model encodes this literally: it has no accept rule that
fires on a single signature.

The distinctness requirement is stated at the level the mechanism can enforce:
the required number of \emph{distinct enrolled approver identities}, with distinct
keys, not merely distinct signature instances. The corresponding machine-checked
properties include, in the state-machine models, SelfApprovalImpossible together
with the properties that no enrolled identity fills two quorum slots and no key
fills two quorum slots, and, in the symbolic quorum model, that a 2-of-2 commitment
requires signatures under two distinct enrolled identities over the same action
bound to the same initiator, that the initiator cannot self-approve, and that no
single signer fills the quorum (Section 7.3). Together these rule out one enrolled
identity filling two slots and one credential being presented twice. They do not
prove that the identities correspond to independent natural persons; that depends
on the directory's identity-proofing controls. Where a policy requires ordered
approval, the chain is checked for acyclicity and linearity.

One residual vector in the multi-approver setting survives all signature checking
and is stated plainly. Since each approver signs their own context, a malicious
orchestrator can show different approvers different optional-member content (for
example, different initiator attestations), and every individual signature remains
valid. The specification therefore requires cross-context consistency: an optional
member present in any context of a receipt must be present, byte-identical in
canonical form, in every context, so every approver demonstrably signed the same
stated reason and attribution. Verifiers implementing the members flag violations;
the rule is a conformance check layered above signature validity, not a change to
it.

Delegated approval authority is constrained rather than prohibited: a delegation
record must appear in the receipt's key proofs, a delegate's effective scope is
the intersection of the delegation grant and the principal's authority at signing
time, chains are bounded in depth, and every link is independently signed.

\section{Host-Record Binding}

Receipts do not live alone. The agent-action record formats now in development
(post-execution capsules, pre-execution permits, audit-trail records, provenance
graphs, intent chains) almost uniformly make the same, sound scoping decision:
human authorization is somebody else's format. The result is a set of
reserved-but-empty slots: an approver disposition whose authority is opaque, an
authority-context stub, a human-override field with a privacy carve-out where the
human should be, a ``signed grant'' whose format is an unassigned work item. A
dozen slots with a dozen ad-hoc answers would be worse than none. The binding
profile is composition \emph{around} the analyzed core of Sections 3 and 4, not a
second protocol: a bound artifact verifies under its own specification, and the
receipt whose bytes it references is the one the symbolic models are about.

The binding profile defines, host-agnostically, what filling such a slot means.
Evidence is bound either \textbf{by reference} (a member carrying the SHA-256
digest of the authorization artifact's canonical bytes, an artifact-type token,
and an optional locator hint) or \textbf{embedded} (a compact self-describing
claim carrying the action digest, mode, named approvals with each approver's own
signature object, and optional quorum semantics). Five requirements make the
binding mean the same thing in every host:

\begin{itemize}[leftmargin=*]
\item \textbf{B1, digest grounding.} A binding is credited only against artifact
bytes: the reference form by digest equality, the embedded form by verifying the
contained signatures. A host field that merely asserts ``a human approved,'' with
neither, must not be treated as human-authorization evidence.
\item \textbf{B2, action agreement.} When the host record binds an action, the
authorization artifact's action binding must agree with it. An artifact
authorizing a different action invalidates the binding; it does not merely weaken
it. This defeats splicing a genuine artifact from one action into another's
record, the confused-deputy case.
\item \textbf{B3, verified versus accepted.} Verifying a binding (digests and
signatures hold) is distinct from accepting it (the relying party trusts the
artifact's issuer via out-of-band pinned key material). A verifier must report the
two separately, and a binding from an unpinned issuer must not be accepted; a
self-issued artifact cannot be self-accepted.
\item \textbf{B4, fail-closed absence.} The absence of a binding is the absence of
evidence, never a default. A relying party whose policy requires human
authorization treats an unbound or unresolvable binding as insufficient. Absence
becomes positive evidence only through a signed \emph{observed-absence} statement:
an attestation that the verifier performed a defined search against defined
sources at a stated time and found none, attesting the search and its emptiness,
never a universal negative.
\item \textbf{B5, consistency.} When a host record carries both forms, the
embedded claim's canonical bytes must hash to the reference's digest; a mismatch
invalidates both. The two forms may not tell two stories.
\end{itemize}

For transparency-logged signed statements (the SCITT host family, RFC 9943), the
profile fixes the digest choice per host-statement profile: either a digest over
the receipt payload's canonical bytes, for offline composition, or a digest over
the COSE\_Sign1 bytes of the receipt expressed as a signed statement, recommended
when the receipt is registered in a transparency service so the reference also
resolves to a logged, inclusion-proofed entry. The reference form additionally
serves disclosure minimization: the approver identifier travels only in the artifact,
which can be withheld until a relying party with a need to know requires it, at the
cost, by B4, of the evidence not counting until produced.

Two boundary rules complete the profile. A host record must not extend an
artifact's validity by carrying it; replay and one-time-use semantics belong to
the artifact and its enforcement point. And the profile defines no authorization
format of its own: a bound artifact verifies under its own specification.

\section{Evidence Sufficiency: the Action Evidence Graph}

\subsection{The relying party's question}

The standards landscape now produces many signed artifacts about one agent action:
workload identity credentials, delegation and grant tokens, call-chain transaction
tokens, runtime attestation results, pre-execution permits, action-bound
human-approval receipts, post-execution records, transparency-log inclusion
receipts. Each answers its own question. None answers the relying party's: given
all of these artifacts, is this action's evidence sufficient \emph{for this
reliance purpose}? A bank releasing a wire, an insurer paying a claim, and a
regulator auditing an agent need different sufficiency bars over the same
artifacts, and today each such decision is ad hoc, non-reproducible, and leaves no
evidence of its own. As with the binding profile, this layer composes \emph{around}
the analyzed core: it consumes the receipt as one node among several and inherits,
without strengthening, the trust model of each artifact it references.

\subsection{Content-addressed graphs, edges as claims}

An Action Evidence Graph is a portable JSON document carrying the action digest,
nodes, and edges. Each node's identifier is the SHA-256 digest of the referenced
artifact's JCS-canonical bytes, with a type token and, optionally, the artifact
inline; a verifier rejects an inline artifact whose canonical digest does not
equal its node id. Because nodes are references, a presenter can disclose the
\emph{shape} of its evidence without the contents; an undisclosed node contributes
nothing to sufficiency, and a required undisclosed node fails closed.

Edges are presenter claims, and this framing carries the security weight of the
layer. An edge (for example, ``this permit \emph{permits} that authorization
receipt'') is credited only if the claimed binding is present in the source
artifact's own canonical bytes, at minimum containment of the target's digest. A
claimed edge that the bytes do not back \emph{poisons} the evaluation: the graph
has asserted something its evidence does not support, and the verdict must not be
admissible.

Graph identity is disclosure-independent: the graph digest is computed over the
sorted (id, type) node pairs, the sorted edges, and the action digest, and never
over inline artifact bytes, so two parties holding different disclosures of the
same graph agree on what graph they are discussing.

\subsection{Policy replay and the closed verdict set}

The sufficiency bar is an evidence policy supplied by the relying party. The graph
document has no policy member, by construction: a presenter choosing its own
sufficiency bar is the confused-deputy failure of evidence systems. A policy names
its reliance purpose, a boolean requirement expression over node types, per-type
freshness bounds, per-type revocation requirements, and required edges that must be
present and byte-backed.

Evaluation is deterministic and offline: given the same graph, policy, and
evaluation time, any implementation reaches the same verdict and the same replay
digest. The verdict is one of a \emph{closed} set of five: \code{admissible},
\code{missing\_evidence}, \code{stale}, \code{conflicted}, \code{unverifiable},
with precedence \textbf{unverifiable over conflicted over stale over
missing\_evidence over admissible}, so no failure ever degrades toward
admissibility. Machine-readable reason codes accompany the verdict; the reason
vocabulary is open, the verdict set is not.

The failure semantics encode a deliberate distinction between absence and
deception. A node referenced but not disclosed contributes an unverified fact; if
the policy required that type, the verdict is \code{missing\_evidence}: the
evidence may exist but was not presented, and nothing about the graph lied. A
claimed edge the source bytes do not back is a lie about the evidence and forces
\code{unverifiable} regardless of what else verified. Absence degrades to ``not
enough''; deception degrades to ``not trustable.''

The replay digest is the canonical digest of exactly three inputs: the policy as
supplied, the ordered list of normalized evidence facts (per node: type, label,
verified flag, attested action digest, outcome if any, revocation flag, age,
derived staleness, and whether a revocation check was required), and the evaluation
time, bound to the graph digest. Raw artifact bytes, network state, and clock reads
are excluded by construction; two evaluators that disagree on a replay digest are,
by definition, evaluating different inputs.

\subsection{Reliance results and policy packs}

The verdict may itself be issued as a signed artifact carrying the verdict,
reasons, action digest, graph digest, policy digest and identifier, reliance
purpose, replay digest, and evaluation time, signed by the relying party's
verifier key. The signature adds accountability, never authority: any third party
can recompute the replay digest rather than trust the signer, and the
verified/accepted distinction of B3 applies to the result itself. Signed reliance
results make reliance decisions auditable: an institution can prove, later, which
policy it applied to which evidence before it acted.

Six policy packs profile the mechanism for concrete irreversible action classes:
wire transfer, vendor bank-account change (a distinct-human quorum), credential
rotation, production data deletion, regulated trade execution (post hoc: the
execution record must provably reference its authorization and be
transparency-logged), and healthcare record export. A pack is a starting point the
relying party adopts and owns, pinning its own trust anchors and tightening
freshness to its risk appetite; no pack waives fail-closed behavior.

A verdict of \code{admissible} is evidence that a bundle met a stated policy at a
stated time. It is not adjudication, it does not establish that the action was
correct or safe, and it confers no authority beyond what the underlying artifacts
carry. The token is a protocol token naming sufficiency under the stated policy; it
carries no legal meaning.

\section{Formal Analysis}

The central formal contribution of this work is a symbolic-security analysis, under
a Dolev-Yao attacker, of the receipt core, a 2-of-2 quorum instance, and two
composed acceptance models spanning authority, registry view, reliance
configuration, challenge consumption, and fleet execution. We present that
analysis first (Sections 7.1 through 7.4), then the
complementary state-machine and relational model checking that surrounds it
(Section 7.5), and finally the discipline that keeps \emph{specified} separate from
\emph{verified} (Section 7.6).

\subsection{The symbolic model and its attacker}

The four Tamarin models are the EP models in which ``the signature verifies'' is a
\emph{derived} fact rather than an axiom. Signatures are terms in Tamarin's
standard \code{signing} equational theory, \code{verify(sign(m, k), m, pk(k)) =
true}, and the adversary is the standard Dolev-Yao network attacker: full control
of the network, sees every message, may apply all function symbols, may extract a
device key through an explicit compromise rule, and may additionally request an
honest human to sign arbitrary actions of the attacker's own choosing. The
attacker's ability to collect honest signatures over any actions it likes is what
makes the no-replay-across-actions results meaningful: the prover must rule out
repurposing any of those honest signatures onto a different action.

The interesting part is the adversary, not the tool. A symbolic analysis usually
lets honest principals sign only what a correct protocol run dictates. Here the
party running the protocol is untrusted for evidence purposes, and it can drive an
honest human to sign any action it chooses, which is the realistic shape of an
operator that controls both the signing client and the approval workflow and later
has an interest in the record. The models ask whether, against that
operator-untrusted, signature-harvesting adversary, an accepted receipt still
implies a prior user-verification-gated signature under the pinned enrolled key
over exactly the accepted action, whether a satisfied quorum still requires two
distinct enrolled identities rather than one signature counted twice, and whether
the one-time-consumption and initiator bindings are what carry those properties.
The models answer yes under their stated abstractions, and dropping either binding
produces a counterexample. They do not prove natural-person identity, independent
human judgment, WebAuthn internals, or downstream exactly-once effects. The
primitives are standard; the contribution is the security argument about them
under this adversary, which is the adversary that the evidence-versus-testimony
problem of Section 1 actually presents.

The models were checked with tamarin-prover 1.10.0 (Maude 3.4), all
well-formedness checks passing. \code{ep\_receipt\_core.spthy} analyzes a single
user-verified signature, its acceptance, and one-time consumption.
\code{ep\_quorum\_core.spthy} layers a 2-of-2 quorum on top of the same
user-verification-gated signature discipline. \code{ep\_reliance\_composed.spthy}
re-derives an abstract 2-of-2 approval path together with challenge, action
identifier, profile and audience, initiator, authority, exact registry view,
revocation, issuer pinning, consumption, and execution.
\code{ep\_six\_claim\_composed.spthy} focuses on Class-A downgrade refusal,
denial non-authorization, scoped authority, relying-party-pinned profiles,
registered challenge consumption, and canonical action-keyed fleet execution.
The models do not cover full WebAuthn internals, directory publication, or
transparency-log mechanics; Section 7.4 states the exact composition and
exclusions.

\subsection{Receipt core: verified lemmas and the consumption falsification}

The core model maps to the receipt construction of Section 3. A human approver
holds a device-bound key whose public half is published (so the attacker always
knows it). User verification is modeled structurally: the signing rule can consume
only a linear fact that the user-verification step alone produces, so no trace
contains a human signature without a prior user-verification event over exactly
that action and nonce. The relying party pins the human's key out of band
(abstracting the approver directory as a single step) and accepts a receipt only if
the signature verifies under the pinned key over the exact action received.
One-time consumption is modeled as a per-(human, nonce) restriction on a
consumption-checked accept rule, sitting beside an unchecked accept rule that omits
it.

The machine-checked result, verbatim from the run, was:

\begin{tool}
executable_honest_receipt (exists-trace): verified (8 steps)
core_authenticity_uv_gated (all-traces): verified (12 steps)
no_replay_across_actions (all-traces): verified (12 steps)
injective_acceptance_with_consumption (all-traces): verified (6 steps)
unchecked_acceptance_is_injective (all-traces): falsified - found trace (10 steps)
\end{tool}

What each verified lemma establishes:

\begin{itemize}[leftmargin=*]
\item \code{executable\_honest\_receipt} (verified). A sanity trace exists: an
honest, user-verified, consumption-checked receipt is accepted with no key
compromise, so the model is not vacuous.
\item \code{core\_authenticity\_uv\_gated} (verified, authenticity gated on user
verification). If any relying party accepts a receipt naming human H for action a
with nonce n, and H's key was not compromised before acceptance, then H performed a
user-verification event and then a signature over exactly (a, n), in that order,
before the acceptance.
\item \code{no\_replay\_across\_actions} (verified, no replay across actions). No
trace exists in which a receipt for action a is accepted while the uncompromised
human never signed exactly (a, n). The attacker holds honest signatures over
arbitrary other actions of its choice and still cannot transplant any of them onto
a different action.
\item \code{injective\_acceptance\_with\_consumption} (verified, injective
acceptance under one-time consumption). With the one-time consumption check, each
consumption-checked acceptance corresponds to a preceding honest signature and no
second consumption-checked acceptance of the same (H, a, n) exists anywhere.
\end{itemize}

The fifth lemma, \code{unchecked\_acceptance\_is\_injective}, is falsified \emph{by
design}, and this is a strength rather than a defect. It asserts injective
acceptance \emph{without} a consumption check. Tamarin falsifies it and produces
the counterexample: the attacker takes the one honest receipt broadcast by the
signing rule and delivers it to an accept point twice. No forgery and no key reveal
occur in the trace; it is a pure same-receipt replay. The verified authenticity
lemma forces every acceptance to be backed by a signature over exactly (a, n), and
the fresh, unique nonce admits at most one such signature, so any double acceptance
is necessarily the single honest receipt accepted twice. This is precisely the
failure the specification's one-time consumption record exists to prevent, and the
model shows mechanically that adding the consumption restriction restores
injectivity. The consumption check is load-bearing, not defense-in-depth, and the
prover demonstrates exactly that.

\subsection{Quorum: verified lemmas and the initiator-binding falsification}

The quorum model maps to Section 4. It enrolls two distinct approver identities,
each with its own device-bound key; a restriction enforces that a single name
cannot be enrolled twice, so the two quorum slots are two distinct names. An
initiator identity, chosen by the attacker, proposes the action, and the initiator
is bound \emph{into} each approver's signed authorization context, abstracted as
\code{<'ep\_signoff\_v1', h(action), initiator, nonce>}. Each approver signs its
own context, user-verification-gated as in the core model. Commitment fires only
when two signatures verify under two distinct pinned approver keys over the same
action and the same initiator, and neither approver is the initiator; there is no
accept rule that fires on a single signature.

The machine-checked result, verbatim from the run, was:

\begin{tool}
executable_quorum (exists-trace): verified (12 steps)
quorum_requires_two_distinct_uv_gated_signatures (all-traces): verified (20 steps)
initiator_cannot_self_approve (all-traces): verified (4 steps)
no_single_signer_fills_quorum (all-traces): verified (4 steps)
commit_requires_signature_over_that_action (all-traces): verified (7 steps)
\end{tool}

What each verified lemma establishes:

\begin{itemize}[leftmargin=*]
\item \code{executable\_quorum} (verified). A 2-of-2 quorum can commit with two
distinct honest approvers, both user-verified, neither the initiator, no key
compromise, so the distinctness and self-approval restrictions do not make
commitment unreachable.
\item \code{quorum\_requires\_two\_distinct\_uv\_gated\_signatures} (verified).
Whenever the executor commits action a citing approvers H1 and H2, and neither
approver key was compromised before the commit, then H1 and H2 are distinct and
each performed a user-verification event followed by a signature over exactly
action a, all before the commit. A single signature, two signatures from one
identity, or two signatures over different actions can never commit.
\item \code{initiator\_cannot\_self\_approve} (verified). No trace commits an
action with the initiator itself appearing as either quorum approver. Because the
initiator identity is attacker-chosen, this also rules out an attacker naming a
compliant approver as initiator to have it count against itself.
\item \code{no\_single\_signer\_fills\_quorum} (verified). The two committing
approvers are never the same identity; distinct pinned keys entail distinct
enrolled identities. This is the no-key-fills-two-slots property at the identity
level.
\item \code{commit\_requires\_signature\_over\_that\_action} (verified). No trace
commits action a while an uncompromised named approver never signed exactly a.
Because the attacker can obtain honest signoffs over arbitrary other actions, this
rules out transplanting an approver's signature over any other action onto a
committed a.
\end{itemize}

As with the consumption check, an earlier revision of this model was falsified, and
the record is kept because it changed the construction. In that first revision the
signed authorization context did \emph{not} bind the initiator identity; it was
\code{<'ep\_signoff\_v1', h(action), nonce>}. Under that revision Tamarin falsified
both \code{quorum\_requires\_two\_distinct\_uv\_gated\_signatures} and
\code{commit\_requires\_signature\_over\_that\_action}: two honest approvers signed
action a while the request carried one initiator label, and the executor committed
the same a under a \emph{different} initiator label, because nothing tied the
approver signatures to the committing initiator. Separation of duties was only an
executor-local check, not a property of the signatures. The fix is spec-faithful
and not a lemma weakening: the initiator identity is now bound inside what each
approver signs, and with that binding all five lemmas verify. The falsification
identified a load-bearing binding, and the model now shows that binding is what
carries the property. This mirrors the core model's consumption result exactly: a
check the design depends on is shown to be depended upon.

\subsection{Composed acceptance paths and the exact scope}

The third model, \code{ep\_reliance\_composed.spthy}, puts the following in one
symbolic trace: a signed reliance challenge; a computed action identifier binding
family and material fields; exact profile, audience, and initiator binding; two
distinct user-verification-gated approvals; scoped authority under an exact pinned
registry epoch and head; revocation state; a pinned receipt issuer; one-time
consumption; and execution. Its machine-checked summary is:

\begin{tool}
executable_composed_reliance (exists-trace): verified (19 steps)
execution_requires_full_composition (all-traces): verified (97 steps)
caid_binds_family_and_material (all-traces): verified (2 steps)
initiator_cannot_self_approve (all-traces): verified (4 steps)
no_single_signer_fills_quorum (all-traces): verified (2 steps)
no_issuer_laundering (all-traces): verified (781 steps)
strict_registry_view_is_exact (all-traces): verified (25 steps)
no_cross_action_profile_or_audience_replay (all-traces): verified (37 steps)
execution_has_honest_approvals_or_prior_compromise (all-traces): verified (170 steps)
injective_execution_with_consumption (all-traces): verified (2 steps)
unchecked_composition_is_injective (all-traces): falsified - found trace (31 steps)
unchecked_registry_view_is_current (all-traces): falsified - found trace (20 steps)
\end{tool}

The two falsified lemmas are deliberately weakened comparisons. Without consumption,
the prover produces a same-receipt replay. Without exact registry-head binding, it
produces acceptance under a stale or equivocating authority view. The strict paths
verify, demonstrating that both checks are load-bearing.

The fourth model, \code{ep\_six\_claim\_composed.spthy}, narrows the second
composed path to six enforcement claims and two companion sanity properties. Its
machine-checked summary is:

\begin{tool}
executable_six_claim_composition (exists-trace): verified (10 steps)
signed_denial_remains_verifiable_evidence (exists-trace): verified (7 steps)
class_a_downgrade_refused (all-traces): verified (22 steps)
signed_denial_cannot_authorize (all-traces): verified (2 steps)
scoped_authority_is_pinned (all-traces): verified (49 steps)
reliance_requires_pinned_profile (all-traces): verified (12 steps)
evidence_challenge_is_registered_and_consumed (all-traces): verified (41 steps)
fresh_challenge_registration_is_unique (all-traces): verified (2 steps)
aec_execution_is_action_keyed_and_fleet_fail_closed (all-traces): verified (13 steps)
action_reservation_failure_is_fail_closed (all-traces): verified (3 steps)
unchecked_presenter_class_is_pinned (all-traces): falsified - found trace (13 steps)
unchecked_signed_denial_cannot_authorize (all-traces): falsified - found trace (13 steps)
unchecked_authority_scope_is_pinned (all-traces): falsified - found trace (15 steps)
unchecked_reliance_profile_is_pinned (all-traces): falsified - found trace (13 steps)
unchecked_unregistered_challenge_is_registered (all-traces): falsified - found trace (13 steps)
unchecked_presenter_execution_key_is_canonical (all-traces): falsified - found trace (13 steps)
\end{tool}

Each unsafe comparison removes one named restriction and produces the expected
attack trace. The challenge result is intentionally narrow: it proves
registration-before-reliance and one-time symbolic consumption, not HTTP
semantics or database durability. Those are separately exercised by bounded
state exploration and executable restart/concurrency tests.

The honest scope statement matters as much as the lemma list. Across the four
models, the analysis does \textbf{not} model, and therefore proves nothing about:

\begin{itemize}[leftmargin=*]
\item \textbf{Approver-directory publication}, \textbf{Merkle-log mechanics},
\textbf{checkpoints}, \textbf{inclusion proofs}, or transparency-log completeness.
The composed model binds an exact pinned registry view as a symbolic value; it does
not prove how that view was published or discovered.
\item \textbf{WebAuthn internals}, authenticator hardware, or clientDataJSON
parsing. User verification is an atomic gating event: the models assume the
authenticator enforces user verification before releasing a signature. They do not
prove WebAuthn.
\item \textbf{Arbitrary k-of-n.} The quorum path fixes the 2-of-2 instance, the
smallest case that exhibits distinctness and self-approval. The parametric k-of-n
statement is not proven.
\item \textbf{Collusion, coercion, and one-human-many-identities.} Separation of
duties defeats unilateral self-approval and guarantees pairwise-distinct signing
\emph{identities}; it does not establish that two identities are two independent
wills.
\item \textbf{Canonicalization, amount arithmetic, policy authorship, or wall-clock
freshness.} Canonical encoding is symbolic and injective; amount and policy
semantics are abstract; no clock model is present. Those surfaces require parser,
arithmetic, policy-governance, and freshness evidence outside this proof.
\item \textbf{Computational or algorithm-specific security, and downstream
exactly-once effects.} \code{sign} is the abstract Dolev-Yao signature with perfect
cryptography; no reduction is offered for ES256, Ed25519, or ML-DSA. Consumption is
proven within the modeled acceptance path, not across an arbitrary downstream
side-effect system.
\end{itemize}

Symbolic verification does not prove the current validity or business correctness
of any real receipt; it proves properties of the construction under these
abstractions.

\subsection{Complementary state-machine and relational model checking}

The symbolic analysis above is about the protocol under a network attacker.
Separately, and complementarily, the surrounding state machine is exhaustively
model-checked at the state-machine level, where signature validity is
\emph{assumed} and the object of study is the reachable state space rather than an
active attacker.

The core state machine is a TLA+ model checked with TLC 2.19. The checked
configuration explores the handshake lifecycle, signoff flow, delegation,
revocation and consumption races, and identity-continuity extensions, generating
413,137 states (45,342 distinct) with no counterexample against 26 invariants.
These include the properties an evidence artifact stands on: an authorization is
never consumed twice; a consumed authorization is never subsequently revoked, and a
revoked one never consumed; the lifecycle is acyclic and terminal states are
inescapable; nonces are unique; delegation chains are acyclic and delegates cannot
exceed their principals' authority; direct-write bypass of the state machine is
impossible; concurrent revoke and consume interleavings serialize; and a denied
action can never become approved. Model checking earned its keep during
development: TLC surfaced four real specification bugs, all fixed before the
properties passed. The checked configuration is a single handshake and a single
claim, which covers all per-handshake and per-claim safety properties exhaustively
within those bounds; it is not a general proof for unbounded concurrent instances,
and the two-handshake configuration was bounded-tested (millions of distinct
states, no counterexample) rather than exhausted. Exhaustively proving the
composed, concurrent case is exactly the kind of question better suited to the
symbolic prover than to TLC brute force, which is one reason the Tamarin models
exist.

The quorum construction, capability delegation, and the protocol's relational
structure are separately checked in four Alloy models (6.2.0, SAT4J), with
32 assertions holding without counterexample. One model proves fifteen assertions over the
receipt and signoff relations (no double consumption, revoked-never-consumed,
consumed-was-verified, terminal-state integrity, write-path exclusivity, delegation
scope containment and acyclicity, signoff binding integrity and consume-once, and
full-chain integrity). A federation model proves seven assertions about the
cross-operator verification path (an accepted receipt is signed by an advertised
key over an untampered payload, a tampered or unadvertised-key or revoked receipt
is never accepted, historical keys still verify, no trust laundering, and
observer-independent portability). A quorum model proves six assertions covering
requester exclusion, distinct humans and keys, the two-person rule, and ordered-chain
acyclicity and linearity. A capability-delegation model proves four assertions
covering acyclicity, unique delegation identifiers, leaf ancestry, and
non-increasing authority. These relational checks treat Ed25519
unforgeability as an abstract fact; they are structural, not cryptographic, which
is precisely the gap the Tamarin models close.

The division of labor is deliberate. Alloy and TLA+ answer ``does the state machine
keep its invariants across the reachable space, assuming signatures behave'';
Tamarin answers ``do the security properties survive an attacker who controls the
network and can solicit honest signatures, with signature validity derived rather
than assumed.'' The three are complementary, and the paper does not present the
state-machine model checking as a substitute for the protocol proof or the reverse.

\subsection{Specified versus verified, and residual normative gaps}

The distinction the project's own status tracking draws is the right one to repeat
here: \emph{specified} means a property is written down as a theorem;
\emph{verified} means a model checker has run and found no counterexample under the
model's assumptions. The two are conflated in this field often enough that keeping
them separate is itself a contribution to hygiene.

Two consequences follow and are stated normatively. First, several parts of the
specification are specified but not covered by any model: the WebAuthn challenge
binding beyond the atomic user-verification assumption, the approver directory
protocol, log checkpoints, and the optional initiator-attestation member are
specified, not proven, and extending the models to them is tracked work. Second, a
few normative mechanisms are ahead of the reference implementation and not yet
exercised by it or by conformance vectors, notably the operator-signed-directory
assurance downgrade, delegation records in the receipt's key proofs together with
the delegate-cannot-exceed-principal check, and enforcement-class emission; on these
the specification is authoritative and the reference implementation is incomplete,
not the reverse. Implementations must not represent any model as covering
deployment topologies or components it does not model.

\section{Implementation and Conformance}

A reference implementation of the receipt, quorum, evidence-chain, binding, and evidence-graph layers is published under Apache-2.0. The conformance battery (21 suites, 331 vectors) is self-contained: each vector carries its own public key and document, requiring no server or shared state. The conformance battery of 21 suites comprising 331 vectors is adversarial (reject vectors each pin one invariant: tampered bytes, wrong key, replay, malformed signature, broken Merkle anchor) and deterministic (regenerable byte-identically from fixed seeds). Further suites cover quorum
semantics, device signoffs, revocation timing, trusted-time attestation, the full
trust receipt with Merkle inclusion and signed checkpoint (including a
transparency-statement digest profile), provenance chains, evidence records, a
canonicalization-malleability battery, the offline-verification boundary, currency
(authentic-as-of-commit versus current validity), initiator-attestation field
validation, sparse-Merkle one-time-consumption transitions, witness cosignatures
over a checkpoint head, and an independent RFC 3161 timestamp-proof profile.

The implementation status is stated precisely, because it is the most tempting
place to overclaim. The JavaScript, Python, and Go reference verifiers live in
\emph{one repository}: they are same-team ports whose agreement across all 331
current vectors is cross-language consistency evidence, not three independent
implementations. A historical Ed25519-signed
\code{EP-EXTERNAL-VERIFICATION-STATEMENT-v1} from COSA / J Diesel NY records
external reproduction of 158 vectors against a pinned commit using this project's
own \code{emilia-verify} package; it remains exactly that historical reproduction
result.

Separately, a Rust verifier authored by J Diesel NY is rebuilt by CI from immutable
public source commit \hexid{7faba36010e7590727bebbc5b9dcceee60539b9b} and tree
\hexid{0553c5fa0a5c4b566703e0d8ef9864dc33e5176d}. The evaluator runs it over the
pinned 16-suite/164-vector clean-room bundle and a pinned hostility campaign of 353
structured attacks plus 6 raw-parser refusals, with zero divergences. The aggregate
conformance case binds both campaigns to the same runner bytes. This is external
interoperability and parser-robustness evidence. Its construction statement is
signed by the implementation organization and predates the pinned source, not by a
distinct attestor over the current source; consequently \code{EP-CONFORMANCE-CASE-v1}
structurally reports zero strict independently attested clean-room acceptances. We
do not upgrade that evidence into a construction-independence claim, and neither
should a reader.

The evidence-graph layer ships with a test suite covering the fail-closed
invariants of Section 6, including disclosure-independent graph identity,
unbacked-edge poisoning, required-edge stripping, freshness degradation, and replay
determinism, together with a complete worked vector (graph, relying-party policy,
and signed reliance result) reproducible from a fixed seed. Deterministic vectors
for the binding profile's checks and for the observed-absence statement are
published alongside.

\section{Related Work}

The receipt is positioned as a composition with adjacent layers, not a replacement
for any of them; each answers a different question, and the composition is the
point.

\textbf{Transparency and logging.} SCITT (RFC 9943) provides an append-only
transparency log with inclusion receipts, and is deliberately agnostic about
\emph{who authorized} a statement, which is precisely the question the
authorization receipt answers. The two compose directly: a receipt can be expressed
as a COSE signed statement and registered in a transparency service, and the
binding profile's statement-digest form (Section 5) is designed for that carriage.
The agent-action statement cluster emerging around SCITT (capsules, permits,
post-execution records, refusal events) makes agent actions transparent, logged, and
policy-checked; those profiles reserve slots for, and do not produce, the enrolled
approver's action-bound approval evidence itself. Receiver-attested logging has the
receiving service sign what it observed, post hoc; pre-execution authorization and
post-execution attestation are complementary halves of a complete record.

\textbf{Workload and agent identity.} WIMSE workload credentials authenticate which
workload is calling and prove key possession; their chartered scope is workload
identity, and they deliberately do not define personal identities. A workload token
is nonetheless carry-capable for the binding profile's reference form, and the two
answers have different trust roots (a workload key versus an enrolled approver key under
organizational authority), which is why neither format can absorb the other. RATS
(RFC 9334) and the Entity Attestation Token (RFC 9711) attest platform or workload
trustworthiness, an orthogonal trust root; the receipt borrows EAT's
verifier-relevant-nonce freshness discipline with a higher floor.

\textbf{Authorization and delegation.} OAuth Rich Authorization Requests (RFC 9396)
and GNAP (RFC 9635) authorize a client's requested scope; step-up authentication
(RFC 9470) can demand fresh human authentication but yields no durable, portable
artifact of the approval. Transaction tokens (draft-ietf-oauth-transaction-tokens)
propagate context across a machine call chain within a trust domain, short-lived and
online-validated; the receipt is the human-approval evidence root from which such a chain
can descend. Delegation-receipt work (draft-nelson-agent-delegation-receipts) binds
a \emph{user's} delegation to an \emph{operator's} instructions, upstream; the
authorization receipt binds an enrolled organizational approver key to an exact action
downstream, with separation-of-duties and quorum semantics the delegation layer does
not formalize, and the two compose by reference. CIBA may serve as the transport by
which an approver is reached; the signoff is what the approver produces on arrival.
Security Event Tokens (RFC 8417) and CAEP convey issuer assertions that an event
occurred, not an enrolled approver key's pre-execution signoff bound to one action.
Per-hop machine-side route and execution receipts govern delegated machine
authority; none binds an enrolled approver key to a specific action under a
human-approval policy, and the receipt's delegation constraints share their
tighten-only scope-narrowing discipline. For long-lived evidence, Evidence Record
Syntax (RFC 4998) supplies the re-timestamping pattern by which receipts remain
verifiable after their original algorithms weaken.

\textbf{Formal methods for authorization protocols.} Symbolic protocol analysis
under a Dolev-Yao attacker with tools such as Tamarin and ProVerif is a mature body
of work; it has been applied to key-exchange and authentication protocols and, more
recently, to WebAuthn and FIDO ceremonies. The contribution here is not the method
but its application to action-bound human-approval evidence: deriving, rather than
assuming, that acceptance of an authorization implies a prior
user-verification-gated signature under a pinned enrolled key over exactly the
accepted action, that quorum requires distinct enrolled identities with no
self-approval, and that the one-time-consumption and initiator
bindings the design depends on are the bindings that carry the properties. The
complementary use of TLA+ for the state machine and Alloy for the relational
structure follows their conventional strengths.

\textbf{Separation from signature security.} The result should not be confused
with a new claim about the underlying signature scheme. Standard signature
security addresses whether an adversary can produce a valid signature on a
message that was never signed; it does not specify whether a verifier accepts
that message twice, accepts it under a different action label, or counts the
same identity in two quorum slots. ABIA makes those verifier-side obligations
explicit. The Tamarin falsifications are consequently separations between
primitive authenticity and protocol authorization: the same ideal signing
algebra is used in the safe and unsafe models, and only the omitted protocol
binding changes.

\section{Limitations}

\subsection{What a receipt does not prove}

A receipt proves that a key enrolled under an approver identifier signed a
canonical context binding one action and that the authorization reached a recorded
terminal state within the relevant consumption domain. It does not prove that a
particular natural person exercised the key: identifier-to-person binding is the
directory and identity-proofing layer's property (Section 3.6), and possession of
the enrolled device plus its user-verification factor is the practical proxy the
artifact rests on. A coerced approver, a shoulder-surfed PIN, or a human who has
enrolled multiple identities all produce receipts that verify. Distinctness
checking operates over enrolled identities; whether two identities are two humans
is an enrollment control, and the symbolic quorum model is explicit that it proves
the distinct-identity property and nothing about independent wills (Section 7.4).
Receipts make insider events attributable to enrolled identities and evidenced,
which raises their cost; they do not make them impossible, and conforming
implementations are prohibited from claiming otherwise.

\subsection{Presentation attacks: narrowed, not solved}

The gravest risk in the construction is that the approver signs context hash H
believing it represents action X when it represents action Y. A signature proves
user presence and an act of approval toward \emph{whatever was rendered};
cryptography cannot prove the rendering was faithful, and the symbolic models do not
attempt to: they assume the human signs the action term they are asked about. The
specification requires escalating mitigations: the signing client must render from
the exact bytes that were hashed, never from a separately supplied description; for
high-value policies, render templates must be registered with the policy and
committed under the policy hash, so display logic is inside what the signature
covers; above an organization-designated threshold, material parameters must
additionally be rendered on a second surface not authored by the orchestrating
operator. Initiator-supplied free text is rendered as untrusted content under a hard
length cap, visually separated from the operator-rendered action.

The residual risk is stated without minimization: absent a trusted display path,
that is, hardware the operator does not author, rendering fidelity is enforced by
controls, audit, and consented mismatch drills, not by mathematics. What the
cryptography does provide is exactness of evidence: the receipt contains the full
Action Object actually signed, so any divergence between what was displayed and what
was executed is detectable after the fact with proof rather than testimony. The
evidence side of the what-you-see-is-what-you-sign problem is closed; the perception
side is narrowed.

A related human-factors limitation is upstream of any attack: a gate that humans
route around protects nothing, and rubber-stamping is the empirical failure mode of
human-in-the-loop controls under volume. The construction is therefore not a general
approval workflow; deployments are directed to scope signoff to genuinely high-risk,
low-frequency actions and to treat signing-latency floors and zero deny rates as
warning signs. The artifact's evidentiary value is only as strong as the attention
of the human at its center.

\subsection{Log completeness is unprovable offline}

Offline verification establishes that a receipt was included in a log tree whose
head the operator signed. It cannot establish, offline, that the tree shown to this
verifier is the tree shown to everyone (equivocation requires gossip or witness
cosigning to detect, which are online activities), nor that the log is complete: the
absence of a receipt for some action is not offline-provable, and becomes evidence
only through the signed observed-absence mechanism of Section 5, which attests a
defined search and its emptiness, never a universal negative. For actions that a
policy gates on signoff at an enforcing deployment class, the absence of any valid
receipt is evidence that the control was bypassed; that property comes from the
gate, not from the log. Revocation currency likewise requires an online consultation
of a current directory head and checkpoint. These are boundaries of the offline
model, stated as such, rather than defects a future version will quietly remove. The
symbolic analysis inherits this boundary directly: it abstracts the log and
directory as out-of-band pinning and proves nothing about split-view or
completeness.

\subsection{Scope of the sufficiency layer}

An evidence-graph verdict is evidence of sufficiency under a stated,
relying-party-owned policy at a stated time. It does not adjudicate disputes, does
not establish business correctness, and inherits, without strengthening or
weakening, the trust model of each artifact class it consumes. The publisher of
these formats is not an auditor, regulator, or insurer; the artifacts are designed
to support the determinations of such parties and never to substitute for them.

\section{Conclusion}

The gap this work addresses is narrow and, we have argued, real: the agent-security
stack now produces abundant signed evidence about what an agent \emph{is}, what it
was \emph{delegated}, what it \emph{did}, and where that was \emph{logged}, while
evidence that an enrolled approver key signed a context binding one exact
irreversible action before terminal consumption typically remains a database row
in the custody of an interested party. The authorization receipt makes that
narrower fact portable and offline-verifiable, with one-time terminal admission
inside a shared atomic consumption domain; the binding profile lets adjacent record
formats carry it without redefining it; and the evidence-graph layer turns ``is
this pile of artifacts enough?'' into a deterministic, replayable,
relying-party-owned computation whose own outcome is evidence.

The design's discipline is subtraction. The operator is removed from the signature
path, then prevented from re-entering through directory authority. The presenter is
removed from the sufficiency decision. Symmetric primitives are removed from the
verification path. Absence is removed from the space of things that can quietly
count as consent. What remains is deliberately small: canonical bytes, approver-held
keys, atomic terminal consumption, and proofs that travel inside the artifact. The
symbolic analysis is what shows that small core actually carries its weight: under
a network attacker with signature validity derived and not assumed, acceptance
implies a prior user-verification-gated signature under the pinned enrolled key
over exactly the accepted action, honest signatures cannot be transplanted,
one-time consumption is what restores injectivity within the modeled domain, and a
satisfied quorum requires distinct enrolled identities over the same action bound
to the same initiator with no self-approval. The two falsification results, the
replay that appears when consumption is dropped and the quorum break that appears
when the initiator is not signed over, are the strongest evidence that these
bindings are load-bearing rather than decorative.

We have been equally deliberate about the boundary of the claims. Two focused
symbolic models cover the receipt core and a 2-of-2 quorum instance; two composed
models analyze challenge, action identifier, approvals, authority, exact pinned
reliance configuration, challenge consumption, and action-keyed fleet execution.
They still exclude WebAuthn internals, directory publication and log mechanics,
arbitrary k-of-n, canonical parsing and amount arithmetic, policy authorship,
clock freshness, collusion, coercion, and downstream effects. The complementary
TLA+ and Alloy models cover the
surrounding state machine with signatures assumed. The cross-language ports are a
consistency check, while the pinned external Rust verifier supplies separately
authored interoperability and hostility evidence; its construction statement is
implementer-signed, so strict independently attested acceptance remains zero.
Rendering fidelity is enforced by controls rather than mathematics. And
verification, everywhere in this design, proves signature, binding, and
log-inclusion integrity as of commit time, never that the authorized action was
correct. Evidence of authorization is not a verdict on conduct. It is the input that
lets whoever must reach such a verdict, an auditor, an insurer, a court, do so from
artifacts rather than testimony, and that is the entire ambition of the design.

\section*{Artifact Availability}

The following are published under the Apache-2.0 license in the EMILIA Protocol
repository: four Tamarin symbolic models (\hexid{ep\_receipt\_core.spthy},
\hexid{ep\_quorum\_core.spthy}, \hexid{ep\_reliance\_composed.spthy}, and
\hexid{ep\_six\_claim\_composed.spthy}) with model
headers, verbatim tool output, falsification history, and re-run instructions; the
TLA+ and Alloy models with checked configurations and a proof-status ledger that
keeps \emph{specified} and \emph{verified} separate; the reference implementation
(receipt, quorum, evidence-chain, host-record binding, and evidence-graph layers,
with JavaScript, Python, and Go same-team ports); the 21-suite / 331-vector
conformance battery with deterministic regeneration scripts; the pinned external
Rust source and evaluator; the 359-case hostility campaign; the aggregate
conformance case; the historical third-party reproduction statement; and the worked
vectors referenced in this paper. The construction is also specified, in engineering
form, as individual Internet-Drafts via the IETF Datatracker
(\hexid{draft-schrock-ep-authorization-receipts} and companion documents);
Internet-Drafts are working documents and should be cited as work in progress.

\begin{thebibliography}{99}

\bibitem{ar07} I. Schrock, ``Authorization Receipts for High-Risk Agent Actions,''
Internet-Draft \code{draft-schrock-ep-authorization-receipts-10}, work in progress,
August 2026.

\bibitem{quorum03} I. Schrock, ``Multi-Party Human Authorization (EP-QUORUM),''
Internet-Draft \code{draft-schrock-ep-quorum-03}, work in progress, July 2026.

\bibitem{aec03} I. Schrock, ``Authorization Evidence Chains: Composing Heterogeneous
Agent-Action Evidence (EP-AEC),'' Internet-Draft
\code{draft-schrock-ep-authorization-evidence-chain-05}, work in progress, August
2026.

\bibitem{er01} I. Schrock, ``Long-Term Evidence Records for Authorization Receipts
(EP-EVIDENCE-RECORD),'' Internet-Draft \code{draft-schrock-ep-evidence-record-01},
work in progress, July 2026.

\bibitem{tamarin} S. Meier, B. Schmidt, C. Cremers, D. Basin, ``The TAMARIN Prover
for the Symbolic Analysis of Security Protocols,'' CAV 2013.

\bibitem{gmr88} S. Goldwasser, S. Micali, R. Rivest, ``A Digital Signature Scheme
Secure Against Adaptive Chosen-Message Attacks,'' SIAM Journal on Computing 17(2),
1988.

\bibitem{dy} D. Dolev, A. C. Yao, ``On the Security of Public Key Protocols,'' IEEE
Transactions on Information Theory, 1983.

\bibitem{tla} L. Lamport, ``Specifying Systems: The TLA+ Language and Tools for
Hardware and Software Engineers,'' Addison-Wesley, 2002.

\bibitem{alloy} D. Jackson, ``Software Abstractions: Logic, Language, and
Analysis,'' MIT Press, revised edition, 2012.

\bibitem{jcs} A. Rundgren, B. Jordan, S. Erdtman, ``JSON Canonicalization Scheme
(JCS),'' RFC 8785, June 2020.

\bibitem{webauthn} W3C, ``Web Authentication: An API for accessing Public Key
Credentials, Level 2,'' April 2021.

\bibitem{scitt} ``An Architecture for Trustworthy and Transparent Digital Supply
Chains (SCITT),'' RFC 9943.

\bibitem{cose} J. Schaad, ``CBOR Object Signing and Encryption (COSE): Structures
and Process,'' RFC 9052.

\bibitem{eat} ``The Entity Attestation Token (EAT),'' RFC 9711.

\bibitem{rats} ``Remote ATtestation procedureS (RATS) Architecture,'' RFC 9334.

\bibitem{gnap} ``Grant Negotiation and Authorization Protocol (GNAP),'' RFC 9635.

\bibitem{rar} ``OAuth 2.0 Rich Authorization Requests,'' RFC 9396.

\bibitem{stepup} ``OAuth 2.0 Step Up Authentication Challenge Protocol,'' RFC 9470.

\bibitem{set} ``Security Event Token (SET),'' RFC 8417.

\bibitem{ers} ``Evidence Record Syntax (ERS),'' RFC 4998.

\bibitem{ciba} OpenID Foundation, ``OpenID Connect Client-Initiated Backchannel
Authentication Flow, Core 1.0,'' September 2021.

\bibitem{delrcpt} R. Nelson, ``Delegation Receipt Protocol for AI Agent
Authorization,'' Internet-Draft \code{draft-nelson-agent-delegation-receipts}, work
in progress, June 2026.

\bibitem{txtoken} ``Transaction Tokens,'' Internet-Draft
\code{draft-ietf-oauth-transaction-tokens}, work in progress.

\end{thebibliography}

\end{document}
